% BHU PYQ -- Environmental Science & Remote Sensing: Environmental Science (2024-25 NEP)
\documentclass[11pt,a4paper]{article}
\usepackage[a4paper,top=2.5cm,bottom=2.5cm,left=2.8cm,right=2.8cm,headheight=14pt]{geometry}
\usepackage{amsmath,amssymb}
\usepackage[T1]{fontenc}
\usepackage[utf8]{inputenc}
\usepackage{lmodern,microtype,fancyhdr,enumitem,hyperref,lastpage,parskip}

\hypersetup{colorlinks=true,urlcolor=black,linkcolor=black,pdftitle={ENVVA21: Environmental Science -- BHU NEP 2024-25}}

\pagestyle{fancy}\fancyhf{}
\renewcommand{\headrulewidth}{0.4pt}\renewcommand{\footrulewidth}{0.4pt}
\fancyhead[L]{\small \textit{Banaras Hindu University}}
\fancyhead[R]{\small \textit{ENVVA21: Environmental Science (2024-25)}}
\fancyfoot[C]{\small Page \thepage\ of \pageref{LastPage}}

\newlist{parts}{enumerate}{2}
\setlist[parts,1]{label=(\alph*),leftmargin=*,itemsep=4pt,topsep=3pt}
\setlist[parts,2]{label=(\roman*),leftmargin=*,itemsep=2pt,topsep=2pt}

\newcommand{\pts}[1]{\hfill \textit{[#1]}}

\begin{document}

\begin{center}
  {\large\bfseries BANARAS HINDU UNIVERSITY}\\[2pt]
  {\normalsize Under Graduate Value Added Course Examination, 2024-25 (NEP)}\\[6pt]
  \rule{0.8\linewidth}{0.5pt}\\[6pt]
  {\large\bfseries ENVIRONMENTAL SCIENCE}\\[4pt]
  {\large\bfseries Paper Code: ENVVA-21 : Environmental Science}\\[4pt]
  {\normalsize Time: 2:30 Hours \hfill Full Marks: 70}\\[2pt]
  {\small REF NO. CE/Conf./D-II/87}
\end{center}
\rule{\linewidth}{0.4pt}
\smallskip

\noindent \textit{Instructions: Attempt all multiple-choice questions. Each question carries 1 mark. Write your Roll No.\ at the top immediately on receipt of this question paper.}

\bigskip

\begin{enumerate}[label=\textbf{\arabic*.},leftmargin=*,itemsep=8pt]

  \item Which of the following is NOT a component of the biosphere?\pts{1}
  \begin{parts}
    \item Lithosphere
    \item Atmosphere
    \item Hydrosphere
    \item Exosphere
  \end{parts}

  \item The main contributors to acid rain are:\pts{1}
  \begin{parts}
    \item Suspended particulate matter and carbon oxides
    \item Sulphur dioxide and nitrogen oxides
    \item Carbon dioxide and carbon monoxide
    \item Nitrogen oxides and carbon oxides
  \end{parts}

  \item The philosophy that encourages collective responsibility to ensure environmental well-being for present and future generations is called:\pts{1}
  \begin{parts}
    \item Environmental ethics
    \item Linear economy
    \item Human development
    \item Industrial revolution
  \end{parts}

  \item Which article of the Indian Constitution states that it is the fundamental duty of every citizen to protect and improve the natural environment?\pts{1}
  \begin{parts}
    \item Article 51A(g)
    \item Article 307
    \item Article 48A
    \item Article 20
  \end{parts}

  \item Which of the following is a keystone species in Sundarbans Biosphere Reserve?\pts{1}
  \begin{parts}
    \item Asiatic Lion
    \item Rhinoceros
    \item Royal Bengal Tiger
    \item Lion tailed Macaque
  \end{parts}

  \item Which mission aimed to free India from open defecation and achieve 100\% scientific management of municipal solid waste?\pts{1}
  \begin{parts}
    \item Swachh Bharat Mission
    \item Unnat Bharat Abhiyan
    \item AMRUT Mission
    \item Ayushman Bharat
  \end{parts}

  \item Which of the following is a wildlife sanctuary in Uttar Pradesh?\pts{1}
  \begin{parts}
    \item Katarniaghat Wildlife Sanctuary
    \item Kaimur Wildlife Sanctuary
    \item Pench Wildlife Sanctuary
    \item Kumbhalgarh Wildlife Sanctuary
  \end{parts}

  \item The greenhouse gas majorly emitted from animal husbandry is:\pts{1}
  \begin{parts}
    \item Methane
    \item Carbon dioxide
    \item Chlorofluorocarbons
    \item Nitrous oxide
  \end{parts}

  \item The primary goal of environmental education is to:\pts{1}
  \begin{parts}
    \item Promote industrialization
    \item Encourage resource exploitation
    \item Increase consumerism
    \item Create awareness about environmental issues
  \end{parts}

  \item In India, electricity generation is majorly from:\pts{1}
  \begin{parts}
    \item Nuclear power plants
    \item Solar power
    \item Coal-fired thermal power plants
    \item Hydro-power
  \end{parts}

  \item A statutory body established in 2003 by the Central Government to implement India's Biological Diversity Act (2002) is:\pts{1}
  \begin{parts}
    \item MoEFCC
    \item CPCB
    \item National Biodiversity Authority (NBA)
    \item NITI Aayog
  \end{parts}

  \item Recently, cheetahs from Namibia and South Africa were translocated to which National Park in Madhya Pradesh?\pts{1}
  \begin{parts}
    \item Panna National Park
    \item Kuno National Park
    \item Kanha National Park
    \item Satpura National Park
  \end{parts}

\end{enumerate}

\end{document}
